\section{Teori Bilangan Dasar}
\label{sec:teori-bilangan}

Kriptografi modern berdiri di atas matematika. Caesar cipher cukup dijelaskan dengan
pergeseran huruf, tetapi RSA menuntut pemahaman tentang bilangan prima dan fungsi
totient, ElGamal dan Diffie--Hellman menuntut pemahaman tentang akar primitif dan
logaritma diskrit, sedangkan AES dan kriptografi kurva eliptik menuntut pemahaman
tentang medan berhingga. Tanpa landasan ini, algoritma kriptografi hanya tampak
sebagai rangkaian langkah yang harus dihafal, bukan sebagai konstruksi yang dapat
dipahami dan dianalisis.

Bab ini menyajikan landasan tersebut dalam dua bagian. Bagian pertama membahas teori
bilangan, yaitu perkakas yang dipakai oleh hampir seluruh algoritma kriptografi yang
akan dibahas pada bab-bab berikutnya. Bagian kedua membahas teori informasi, khususnya
entropi, yang memberi kita cara mengukur seberapa besar informasi yang sebenarnya
dibocorkan sebuah cipherteks.

Sebelum masuk ke rinciannya, Tabel~\ref{tab:peta-matematika} memetakan bidang-bidang
matematika yang menopang kriptografi beserta tempat pembahasannya di dalam buku ini.
Peta ini penting agar pembaca mengetahui bagian mana yang sudah dikuasainya dari mata
kuliah lain dan bagian mana yang benar-benar baru \citep{menezes1996}.

\begin{table}[htbp]
    \centering
    \caption{Peta topik matematika untuk kriptografi dan tempat pembahasannya}
    \label{tab:peta-matematika}
    \small
    \begin{tabularx}{\linewidth}{@{}>{\raggedright\arraybackslash}p{3.1cm}
                                  >{\raggedright\arraybackslash}X
                                  >{\raggedright\arraybackslash}p{3.6cm}@{}}
        \hline
        \textbf{Bidang} & \textbf{Topik} & \textbf{Dibahas di} \\
        \hline
        Teori bilangan dasar & Aritmetika modulo, kekongruenan, PBB, relatif prima,
        invers modulo & Bab ini (ringkasan Matematika Diskrit) \\
        \hline
        Teori bilangan lanjut & Fungsi totient, Teorema Euler, akar primitif,
        logaritma diskrit & Bab ini \\
        \hline
        Probabilitas dan statistik & Sebaran frekuensi, indeks kebetulan, uji
        keacakan & Bab~\ref{chap:kriptografi-klasik} \\
        \hline
        Teori informasi & Entropi, laju bahasa, redundansi & Bab ini \\
        \hline
        Kompleksitas algoritma & Notasi $O$, kelas P dan NP, reduksi &
        Bab~\ref{chap:layanan-keamanan} dan bab-bab kunci publik \\
        \hline
        Aljabar abstrak & Grup, gelanggang, medan, medan berhingga,
        $GF(2^8)$ & Bab~\ref{chap:aes} dan Bab~\ref{chap:ecc} \\
        \hline
    \end{tabularx}
\end{table}

\subsection{Sifat Pembagian pada Bilangan Bulat}

Seluruh aritmetika kriptografi bertumpu pada satu fakta sederhana tentang pembagian.
Jika $a$ adalah bilangan bulat dan $n$ adalah bilangan bulat positif, maka selalu
terdapat tepat satu pasangan bilangan bulat $q$ (hasil bagi) dan $r$ (sisa) sehingga
\begin{equation}
    a = qn + r, \qquad 0 \le r < n.
    \label{eq:pembagian}
\end{equation}
Kenyataan bahwa sisa $r$ selalu dapat dibuat \emph{nonnegatif} dan lebih kecil dari
$n$ inilah yang membuat aritmetika modulo dapat didefinisikan secara tunggal. Nilai
$r$ inilah yang disebut $a \bmod n$.

Kita juga memerlukan notasi keterbagian. Ditulis $n \mid a$ (dibaca ``$n$ membagi
$a$'') bila terdapat bilangan bulat $k$ sehingga $a = kn$; dengan kata lain, bila
$a \bmod n = 0$. Sebagai contoh, $3 \mid 15$ karena $15 = 5 \cdot 3$, sedangkan
$3 \nmid 16$. Notasi ini akan muncul terus-menerus, terutama pada definisi
kekongruenan di Bagian~\ref{sec:kekongruenan}.

Beberapa sifat keterbagian yang sering dipakai: bila $n \mid a$ dan $n \mid b$ maka
$n \mid (a + b)$ dan $n \mid (a - b)$; bila $n \mid a$ maka $n \mid ab$ untuk setiap
bilangan bulat $b$. Semua sifat ini bekerja karena keterbagian pada dasarnya adalah
pernyataan tentang kelipatan \citep{stinson2018}.

\subsection{Aritmetika Modulo}

Aritmetika modulo adalah sistem aritmetika untuk bilangan bulat, di mana angka-angka
``berputar kembali'' setelah mencapai nilai tertentu yang disebut modulus. Jika $a$
dan $n$ adalah bilangan bulat dengan $n > 0$, maka $a \bmod n$ adalah sisa bagi dari
$a$ oleh $n$ sebagaimana didefinisikan pada Persamaan~\ref{eq:pembagian}.

Untuk bilangan positif, perhitungannya langsung:
\begin{itemize}
    \item $23 \bmod 5 = 3$, karena $23 = 4 \times 5 + 3$.
    \item $41 \bmod 9 = 5$, karena $41 = 4 \times 9 + 5$.
    \item $17 \bmod 3 = 2$, karena $17 = 5 \times 3 + 2$.
\end{itemize}

Untuk bilangan negatif, banyak mahasiswa tergelincir. Aturan bakunya adalah: hasil
$a \bmod n$ harus selalu berada pada rentang $0 \le r < n$, sehingga jawaban
\emph{tidak boleh} negatif. Untuk $a < 0$, hasilnya diperoleh dengan mengambil
kelipatan $n$ terbesar yang tidak melebihi $a$, lalu menghitung selisihnya:
\begin{itemize}
    \item $-41 \bmod 9 = 4$, karena $9 \mid -41$ menghasilkan sisa $-5$ menurut
          pembagian biasa, dan $9 + (-5) = 4$. Cara lain: $-41 = (-5)\times 9 + 4$.
    \item $-24 \bmod 11 = 9$, karena $-24 = (-3) \times 11 + 9$.
\end{itemize}
Perhatikan pada contoh kedua bahwa $-3 \times 11 = -33$, dan $-24 - (-33) = 9$.
Memeriksa hasil dengan cara mengalikan kembali seperti ini adalah kebiasaan yang baik.

% Ilustrasi siklus aritmetika modulo: menghitung 23 mod 7.
% Dirujuk dari section-01.tex dengan % Ilustrasi siklus aritmetika modulo: menghitung 23 mod 7.
% Dirujuk dari section-01.tex dengan % Ilustrasi siklus aritmetika modulo: menghitung 23 mod 7.
% Dirujuk dari section-01.tex dengan \input{figures/siklus-modulo}.
\begin{figure}[htbp]
    \centering
    \begin{tikzpicture}[
        titik/.style={circle, draw, fill=blue!8, font=\scriptsize,
                      inner sep=1.2pt, minimum size=5.5mm},
    ]
        \def\radius{1.9}
        % Lingkaran acuan (jam modulo 7)
        \draw[gray!40, thin] (0,0) circle (\radius);
        % Titik-titik sisa bagi 0 sampai 6
        \foreach \i/\sudut in {0/90, 1/38.571, 2/-12.857, 3/-64.286,
                               4/-115.714, 5/-167.143, 6/141.429} {
            \node[titik] (p\i) at (\sudut:\radius) {\i};
        }
        % Arah pencacahan: berputar dari 0 melewati 1 lalu berhenti di 2
        \draw[very thick, red!70, -{Latex[length=2.4mm]}]
            (104:\radius+0.42) arc (104:-6:\radius+0.42);
        \node[font=\small, align=center] at (0,0) {$23 \bmod 7$\\$= 2$};
    \end{tikzpicture}
    \caption{Siklus aritmetika modulo 7: sisa bagi diperoleh dengan \textit{berputar} mengelilingi siklus, sehingga $23 \bmod 7$ berhenti pada titik 2}
    \label{fig:siklus-modulo}
\end{figure}
.
\begin{figure}[htbp]
    \centering
    \begin{tikzpicture}[
        titik/.style={circle, draw, fill=blue!8, font=\scriptsize,
                      inner sep=1.2pt, minimum size=5.5mm},
    ]
        \def\radius{1.9}
        % Lingkaran acuan (jam modulo 7)
        \draw[gray!40, thin] (0,0) circle (\radius);
        % Titik-titik sisa bagi 0 sampai 6
        \foreach \i/\sudut in {0/90, 1/38.571, 2/-12.857, 3/-64.286,
                               4/-115.714, 5/-167.143, 6/141.429} {
            \node[titik] (p\i) at (\sudut:\radius) {\i};
        }
        % Arah pencacahan: berputar dari 0 melewati 1 lalu berhenti di 2
        \draw[very thick, red!70, -{Latex[length=2.4mm]}]
            (104:\radius+0.42) arc (104:-6:\radius+0.42);
        \node[font=\small, align=center] at (0,0) {$23 \bmod 7$\\$= 2$};
    \end{tikzpicture}
    \caption{Siklus aritmetika modulo 7: sisa bagi diperoleh dengan \textit{berputar} mengelilingi siklus, sehingga $23 \bmod 7$ berhenti pada titik 2}
    \label{fig:siklus-modulo}
\end{figure}
.
\begin{figure}[htbp]
    \centering
    \begin{tikzpicture}[
        titik/.style={circle, draw, fill=blue!8, font=\scriptsize,
                      inner sep=1.2pt, minimum size=5.5mm},
    ]
        \def\radius{1.9}
        % Lingkaran acuan (jam modulo 7)
        \draw[gray!40, thin] (0,0) circle (\radius);
        % Titik-titik sisa bagi 0 sampai 6
        \foreach \i/\sudut in {0/90, 1/38.571, 2/-12.857, 3/-64.286,
                               4/-115.714, 5/-167.143, 6/141.429} {
            \node[titik] (p\i) at (\sudut:\radius) {\i};
        }
        % Arah pencacahan: berputar dari 0 melewati 1 lalu berhenti di 2
        \draw[very thick, red!70, -{Latex[length=2.4mm]}]
            (104:\radius+0.42) arc (104:-6:\radius+0.42);
        \node[font=\small, align=center] at (0,0) {$23 \bmod 7$\\$= 2$};
    \end{tikzpicture}
    \caption{Siklus aritmetika modulo 7: sisa bagi diperoleh dengan \textit{berputar} mengelilingi siklus, sehingga $23 \bmod 7$ berhenti pada titik 2}
    \label{fig:siklus-modulo}
\end{figure}


Aritmetika modulo juga memenuhi sifat-sifat yang memudahkan perhitungan, yaitu
penjumlahan dan perkalian dapat dilakukan sebelum maupun sesudah pengurangan modulo:
\begin{equation}
    (a + b) \bmod n = \bigl((a \bmod n) + (b \bmod n)\bigr) \bmod n,
    \label{eq:mod-tambah}
\end{equation}
\begin{equation}
    (a \cdot b) \bmod n = \bigl((a \bmod n) \cdot (b \bmod n)\bigr) \bmod n.
    \label{eq:mod-kali}
\end{equation}
Sifat terakhir sangat berharga dalam praktik. Untuk menghitung $3^{100} \bmod 7$,
kita tidak perlu menghitung $3^{100}$ terlebih dahulu --- bilangan itu memiliki 48
angka --- melainkan cukup mengalikan hasil antara secara berulang sambil sesekali
mengurangi modulo. Teknik inilah yang membuat eksponensiasi modular pada RSA,
Diffie--Hellman, dan ElGamal dapat dijalankan pada komputer biasa.

\subsection{Kekongruenan}
\label{sec:kekongruenan}

Dua bilangan bulat $a$ dan $b$ dikatakan \textbf{kongruen modulo $n$} jika $a$ dan $b$
memiliki sisa bagi yang sama saat dibagi oleh $n$. Hal ini ditulis sebagai
\[ a \equiv b \pmod{n}, \]
yang berlaku jika dan hanya jika $n$ membagi habis selisih antara $a$ dan $b$:
\[ n \mid (a - b). \]
Perhatikan perbedaan notasi yang halus namun penting: $a \bmod n$ adalah sebuah
\emph{bilangan} (sisanya), sedangkan $a \equiv b \pmod{n}$ adalah sebuah
\emph{pernyataan} yang bernilai benar atau salah. Menuliskan $17 \equiv 5$ adalah
salah; menuliskan $17 \equiv 2 \pmod{3}$ adalah benar.

\begin{itemize}
    \item $17 \equiv 2 \pmod{3}$ karena $3$ membagi habis $17 - 2 = 15$.
    \item $-7 \equiv 15 \pmod{11}$ karena $11$ membagi habis $-7 - 15 = -22$.
    \item $23 \equiv 3 \pmod{5}$ karena $5$ membagi habis $23 - 3 = 20$.
\end{itemize}

Kekongruenan memiliki sifat-sifat berikut, dengan $a \equiv b \pmod{n}$ dan
$c \equiv d \pmod{n}$:
\begin{itemize}
    \item $a + c \equiv b + d \pmod{n}$,
    \item $a - c \equiv b - d \pmod{n}$,
    \item $a \cdot c \equiv b \cdot d \pmod{n}$,
    \item $a^k \equiv b^k \pmod{n}$ untuk setiap bilangan bulat positif $k$.
\end{itemize}
Sifat-sifat ini memungkinkan kita mengganti sembarang bilangan dengan bilangan lain
yang kongruen dengannya tanpa mengubah hasil akhir modulo $n$ --- persis seperti
$3/6$ dapat diganti dengan $1/2$ pada aritmetika biasa.

Satu hal yang \emph{tidak} berlaku secara umum adalah pencoretan. Dari
$ac \equiv bc \pmod{n}$ kita tidak boleh langsung menyimpulkan $a \equiv b \pmod{n}$.
Contoh penyangkalnya: $2 \cdot 3 \equiv 4 \cdot 3 \pmod{6}$ bernilai benar (keduanya
$\equiv 0$), tetapi $2 \not\equiv 4 \pmod{6}$. Pencoretan hanya sah bila faktor yang
dicoret relatif prima terhadap modulus, yaitu bila $\text{PBB}(c, n) = 1$.

\subsection{Relatif Prima dan PBB}

Dua bilangan bulat $a$ dan $b$ dikatakan \textbf{relatif prima}
(\textit{coprime}) jika pembagi persekutuan terbesar (PBB) atau \textit{Greatest
Common Divisor} (GCD) dari keduanya adalah $1$:
\[ \text{PBB}(a, b) = 1. \]
PBB sendiri didefinisikan sebagai bilangan bulat positif terbesar yang membagi habis
keduanya. Sebagai contoh, $\text{PBB}(23, 40) = 1$ sehingga 23 dan 40 relatif prima,
sedangkan $\text{PBB}(12, 18) = 6$ sehingga keduanya tidak relatif prima.

Beberapa sifat PBB yang akan kita pakai:
\begin{itemize}
    \item $\text{PBB}(a, 0) = |a|$ untuk setiap $a \neq 0$.
    \item $\text{PBB}(a, b) = \text{PBB}(b, a)$.
    \item $\text{PBB}(a, b) = \text{PBB}(b,\, a \bmod b)$ --- sifat inilah yang
          menjadi dasar Algoritma Euclidean.
\end{itemize}

Konsep relatif prima menentukan keberadaan invers modulo, dan oleh karena itu menjadi
syarat yang harus diperiksa sebelum sebuah kunci dapat dipakai. Sebagaimana akan
dilihat pada Bab~\ref{chap:rsa}, pemilihan eksponen publik $e$ pada RSA pada dasarnya
adalah pemilihan bilangan yang relatif prima terhadap $\varphi(n)$.

\subsection{Algoritma Euclidean}

Untuk menghitung PBB dua bilangan besar, pemfaktoran bukanlah cara yang baik: ia
justru menjadi persoalan sulit yang mendasari keamanan RSA. Yang dipakai adalah
\textbf{Algoritma Euclidean}, yang bekerja hanya dengan pembagian berulang
berdasarkan sifat $\text{PBB}(a, b) = \text{PBB}(b, a \bmod b)$. Algoritmanya
berhenti ketika sisanya nol, dan PBB-nya adalah sisa bukan nol terakhir.

\begin{description}
    \item[Masukan:] dua bilangan bulat $a$ dan $b$ dengan $a > b \ge 0$.
    \item[Keluaran:] $\text{PBB}(a,b)$.
    \item[Langkah:] selama $b \neq 0$, ganti pasangan $(a, b)$ dengan $(b, a \bmod b)$.
          Ketika $b = 0$, jawabannya adalah $a$.
\end{description}

Sebagai contoh sederhana, hitung $\text{PBB}(48, 18)$:
\[
\begin{aligned}
48 &= 2 \cdot 18 + 12\\
18 &= 1 \cdot 12 + 6\\
12 &= 2 \cdot 6 + 0
\end{aligned}
\]
Sisa bukan nol terakhir adalah $6$, jadi $\text{PBB}(48, 18) = 6$. Perhatikan bahwa
proses ini tidak pernah mencari faktor prima sama sekali; ia hanya mengurangi ukuran
bilangan dengan cepat. Jumlah langkahnya sebanding dengan banyaknya angka pada
bilangan yang lebih kecil, sehingga PBB dua bilangan 1024 bit dapat dihitung dalam
sekejap \citep{menezes1996}. Contoh yang lebih panjang, dengan
$\text{PBB}(1970, 1066) = 2$, disajikan pada bagian Contoh Terhitung bab ini.

\subsection{Kombinasi Lanjar}

Algoritma Euclidean dapat diperluas untuk menghasilkan sesuatu yang lebih berguna
daripada sekadar nilai PBB. \textbf{Teorema B\'ezout} menyatakan bahwa PBB dua
bilangan selalu dapat ditulis sebagai kombinasi lanjar (\textit{linear combination})
keduanya dengan koefisien bilangan bulat.

\begin{theorem}[Identitas B\'ezout]
Untuk setiap bilangan bulat $a$ dan $b$ yang tidak keduanya nol, terdapat bilangan
bulat $m$ dan $n$ sehingga
\[ ma + nb = \text{PBB}(a, b). \]
\end{theorem}

Koefisien $m$ dan $n$ diperoleh dengan menjalankan Algoritma Euclidean lalu
melakukan \textbf{substitusi balik} dari langkah terakhir ke langkah pertama. Untuk
$\text{PBB}(48,18)=6$ di atas:
\[
\begin{aligned}
6 &= 18 - 1 \cdot 12\\
  &= 18 - 1 \cdot (48 - 2 \cdot 18) && \text{substitusi } 12 = 48 - 2 \cdot 18\\
  &= 3 \cdot 18 - 1 \cdot 48.
\end{aligned}
\]
Jadi $m = -1$ dan $\text{PBB}(48,18) = 3 \cdot 18 - 1 \cdot 48$, dan pemeriksaan
kembali memberi $54 - 48 = 6$. Nilai koefisien ini tidak tunggal, tetapi algoritma
Euclidean diperluas (\textit{extended Euclidean algorithm}) selalu menghasilkan satu
pasangan yang sah dengan biaya komputasi yang sama dengan algoritma aslinya.

\subsection{Invers Modulo}

Invers modulo dari $a$ modulo $n$ adalah bilangan $x$ sedemikian sehingga
\begin{equation}
    a \cdot x \equiv 1 \pmod{n}.
    \label{eq:invers}
\end{equation}
Notasinya adalah $x = a^{-1} \bmod n$. Jangan tertukar dengan kebalikan biasa: pada
modulo 9, $4^{-1}$ bukan $1/4$ melainkan $7$, karena $4 \cdot 7 = 28 \equiv 1
\pmod 9$.

Invers modulo $a$ modulo $n$ ada jika dan hanya jika $a$ dan $n$ relatif prima, yaitu
$\text{PBB}(a, n) = 1$. Alasannya dapat dilihat langsung dari Identitas B\'ezout:
bila $\text{PBB}(a,n)=1$, maka terdapat $m$ dan $n'$ dengan $ma + n'n = 1$; dengan
mengurangi modulo $n$ suku $n'n$ hilang dan tersisa $ma \equiv 1 \pmod n$, sehingga
$m$ adalah invers yang dicari. Sebaliknya, bila $\text{PBB}(a,n) = d > 1$, maka
$a \cdot x$ selalu kelipatan $d$ sehingga tidak mungkin kongruen dengan $1$ modulo
$n$.

\begin{itemize}
    \item $4^{-1} \bmod 9 \equiv 7 \pmod 9$, karena $4 \cdot 7 = 28 = 3 \cdot 9 + 1$.
    \item $23^{-1} \bmod 10 = -3 \equiv 7 \pmod{10}$, karena $23 \cdot (-3) = -69 =
          (-7)\cdot 10 + 1$.
\end{itemize}
Perhatikan pada contoh kedua bahwa $-3$ dan $7$ sama-sama sah sebagai jawaban karena
keduanya kongruen modulo 10; bentuk yang biasanya dilaporkan adalah bentuk
nonnegatifnya, yaitu $7$.

Kemampuan menghitung invers modulo dengan cepat adalah keterampilan yang akan dipakai
berulang kali. Pada RSA, kunci privat $d$ adalah invers dari $e$ modulo $\varphi(n)$;
pada ElGamal dan kriptografi kurva eliptik, dekripsi menuntut invers modulo orde
kelompok. Pada semuanya, invers dihitung dengan Algoritma Euclidean yang diperluas,
bukan dengan mencoba satu per satu.

% Alur Algoritma Euclidean dan pemakaiannya untuk mencari invers modulo.
% Dirujuk dari section-01.tex dengan % Alur Algoritma Euclidean dan pemakaiannya untuk mencari invers modulo.
% Dirujuk dari section-01.tex dengan % Alur Algoritma Euclidean dan pemakaiannya untuk mencari invers modulo.
% Dirujuk dari section-01.tex dengan \input{figures/alur-euclid-diperluas}.
\begin{figure}[htbp]
    \centering
    \begin{tikzpicture}[
        kotak/.style={draw, rounded corners, align=center, font=\small,
                      inner sep=5pt, minimum width=7.2cm},
        panah/.style={-{Latex[length=2.4mm]}, thick},
    ]
        \node[kotak, fill=gray!10] (masuk) at (0,0)
            {Masukan $a$ dan $n$ dengan $0 < a < n$};
        \node[kotak, fill=orange!12] (bagi) at (0,-1.5)
            {$r \leftarrow a \bmod n$; \quad $a \leftarrow n$; \quad $n \leftarrow r$};
        \node[kotak, fill=orange!12] (cek) at (0,-3.0)
            {Apakah $r = 0$?};
        \node[kotak, fill=green!10] (pbb) at (0,-4.5)
            {$\text{PBB}(a_0, n_0) = a$};
        \node[kotak, fill=blue!6] (invers) at (0,-6.0)
            {Substitusi balik: peroleh $x$ dengan $a_0 x \equiv 1 \pmod{n_0}$};

        \draw[panah] (masuk) -- (bagi);
        \draw[panah] (bagi) -- (cek);
        \draw[panah] (cek) -- node[right, font=\scriptsize, fill=white, inner sep=1pt]
            {ya} (pbb);
        \draw[panah] (pbb) -- node[right, font=\scriptsize, fill=white, inner sep=1pt]
            {PBB $= 1$} (invers);

        % Simpai ulang selama sisa bagi belum nol
        \draw[panah] (cek.east) -- (4.4,-3.0) -- (4.4,-1.5) -- (bagi.east)
            node[midway, above, font=\scriptsize, fill=white, inner sep=1pt] {belum};

        % Cabang kegagalan: PBB bukan 1
        \node[kotak, fill=red!6] (gagal) at (0,-7.5)
            {Jika $\text{PBB}(a_0, n_0) \neq 1$: invers modulo tidak ada};
        \draw[panah] (invers) -- (gagal);
    \end{tikzpicture}
    \caption{Alur pencarian invers modulo: Algoritma Euclidean mengulang pembagian sampai sisa bagi nol untuk memperoleh PBB, lalu substitusi balik menghasilkan $x$ sedemikian sehingga $a x \equiv 1 \pmod{n}$}
    \label{fig:alur-euclid-diperluas}
\end{figure}
.
\begin{figure}[htbp]
    \centering
    \begin{tikzpicture}[
        kotak/.style={draw, rounded corners, align=center, font=\small,
                      inner sep=5pt, minimum width=7.2cm},
        panah/.style={-{Latex[length=2.4mm]}, thick},
    ]
        \node[kotak, fill=gray!10] (masuk) at (0,0)
            {Masukan $a$ dan $n$ dengan $0 < a < n$};
        \node[kotak, fill=orange!12] (bagi) at (0,-1.5)
            {$r \leftarrow a \bmod n$; \quad $a \leftarrow n$; \quad $n \leftarrow r$};
        \node[kotak, fill=orange!12] (cek) at (0,-3.0)
            {Apakah $r = 0$?};
        \node[kotak, fill=green!10] (pbb) at (0,-4.5)
            {$\text{PBB}(a_0, n_0) = a$};
        \node[kotak, fill=blue!6] (invers) at (0,-6.0)
            {Substitusi balik: peroleh $x$ dengan $a_0 x \equiv 1 \pmod{n_0}$};

        \draw[panah] (masuk) -- (bagi);
        \draw[panah] (bagi) -- (cek);
        \draw[panah] (cek) -- node[right, font=\scriptsize, fill=white, inner sep=1pt]
            {ya} (pbb);
        \draw[panah] (pbb) -- node[right, font=\scriptsize, fill=white, inner sep=1pt]
            {PBB $= 1$} (invers);

        % Simpai ulang selama sisa bagi belum nol
        \draw[panah] (cek.east) -- (4.4,-3.0) -- (4.4,-1.5) -- (bagi.east)
            node[midway, above, font=\scriptsize, fill=white, inner sep=1pt] {belum};

        % Cabang kegagalan: PBB bukan 1
        \node[kotak, fill=red!6] (gagal) at (0,-7.5)
            {Jika $\text{PBB}(a_0, n_0) \neq 1$: invers modulo tidak ada};
        \draw[panah] (invers) -- (gagal);
    \end{tikzpicture}
    \caption{Alur pencarian invers modulo: Algoritma Euclidean mengulang pembagian sampai sisa bagi nol untuk memperoleh PBB, lalu substitusi balik menghasilkan $x$ sedemikian sehingga $a x \equiv 1 \pmod{n}$}
    \label{fig:alur-euclid-diperluas}
\end{figure}
.
\begin{figure}[htbp]
    \centering
    \begin{tikzpicture}[
        kotak/.style={draw, rounded corners, align=center, font=\small,
                      inner sep=5pt, minimum width=7.2cm},
        panah/.style={-{Latex[length=2.4mm]}, thick},
    ]
        \node[kotak, fill=gray!10] (masuk) at (0,0)
            {Masukan $a$ dan $n$ dengan $0 < a < n$};
        \node[kotak, fill=orange!12] (bagi) at (0,-1.5)
            {$r \leftarrow a \bmod n$; \quad $a \leftarrow n$; \quad $n \leftarrow r$};
        \node[kotak, fill=orange!12] (cek) at (0,-3.0)
            {Apakah $r = 0$?};
        \node[kotak, fill=green!10] (pbb) at (0,-4.5)
            {$\text{PBB}(a_0, n_0) = a$};
        \node[kotak, fill=blue!6] (invers) at (0,-6.0)
            {Substitusi balik: peroleh $x$ dengan $a_0 x \equiv 1 \pmod{n_0}$};

        \draw[panah] (masuk) -- (bagi);
        \draw[panah] (bagi) -- (cek);
        \draw[panah] (cek) -- node[right, font=\scriptsize, fill=white, inner sep=1pt]
            {ya} (pbb);
        \draw[panah] (pbb) -- node[right, font=\scriptsize, fill=white, inner sep=1pt]
            {PBB $= 1$} (invers);

        % Simpai ulang selama sisa bagi belum nol
        \draw[panah] (cek.east) -- (4.4,-3.0) -- (4.4,-1.5) -- (bagi.east)
            node[midway, above, font=\scriptsize, fill=white, inner sep=1pt] {belum};

        % Cabang kegagalan: PBB bukan 1
        \node[kotak, fill=red!6] (gagal) at (0,-7.5)
            {Jika $\text{PBB}(a_0, n_0) \neq 1$: invers modulo tidak ada};
        \draw[panah] (invers) -- (gagal);
    \end{tikzpicture}
    \caption{Alur pencarian invers modulo: Algoritma Euclidean mengulang pembagian sampai sisa bagi nol untuk memperoleh PBB, lalu substitusi balik menghasilkan $x$ sedemikian sehingga $a x \equiv 1 \pmod{n}$}
    \label{fig:alur-euclid-diperluas}
\end{figure}


\subsection{Bilangan Prima}

\textbf{Bilangan prima} adalah bilangan bulat lebih besar dari 1 yang hanya habis
dibagi oleh 1 dan dirinya sendiri. Bilangan 2, 3, 5, 7, 11, dan 13 adalah prima,
sedangkan 1 bukan prima (inilah konvensi yang dipakai di seluruh buku ini) dan 9
bukan prima karena $9 = 3 \cdot 3$.

Dua fakta tentang bilangan prima menopang seluruh kriptografi kunci publik.
Pertama, \textbf{Teorema Dasar Aritmetika}: setiap bilangan bulat lebih besar dari 1
dapat dinyatakan sebagai hasil kali bilangan-bilangan prima secara tunggal (kecuali
urutan faktornya). Kedua, dari hasil kali dua bilangan prima besar kita tidak memiliki
cara yang praktis untuk mengembalikannya menjadi faktor-faktornya. Inilah asimetri
yang dimanfaatkan RSA: mengalikan dua bilangan prima 1024 bit mudah, memfaktorkan
hasilnya sangat sulit \citep{katz2020}.

Euclid membuktikan lebih dari dua ribu tahun yang lalu bahwa banyaknya bilangan prima
tidak terbatas. Meskipun demikian, bilangan prima menjadi makin jarang seiring
bertambahnya nilai. Sekitar satu dari setiap $\ln x$ bilangan di sekitar $x$ adalah
prima, sehingga untuk memperoleh bilangan prima 1024 bit, pembangkit kunci hanya
perlu menguji beberapa ratus kandidat. Pengujian keprimaannya sendiri tidak dilakukan
dengan mencoba membagi, melainkan dengan uji probabilistik seperti Miller--Rabin yang
dapat menyatakan sebuah bilangan komposit dengan pasti, dan menyatakan bilangan prima
dengan peluang benar yang dapat dibuat sedekat mungkin dengan 1.

\subsection{Fungsi Totient Euler}

\textbf{Fungsi totient Euler}, ditulis $\varphi(n)$, menyatakan banyaknya bilangan
bulat pada rentang $1 \le k \le n$ yang relatif prima terhadap $n$. Sebagai contoh,
untuk $n = 20$ bilangan yang relatif prima terhadap 20 adalah
\[ 1,\ 3,\ 7,\ 9,\ 11,\ 13,\ 17,\ 19, \]
seluruhnya delapan buah, sehingga $\varphi(20) = 8$. Bilangan genap, kelipatan 5, dan
kelipatan 10 tidak dihitung karena semuanya memiliki faktor persekutuan dengan 20.

Untuk keperluan kriptografi, tiga sifat berikut sudah mencukupi:
\begin{itemize}
    \item $\varphi(1) = 1$ dan $\varphi(p) = p - 1$ untuk setiap bilangan prima $p$.
          Untuk $p = 7$ diperoleh $\varphi(7) = 6$, sesuai dengan bilangan
          $1,2,3,4,5,6$ yang semuanya relatif prima terhadap 7.
    \item $\varphi(p^k) = p^k - p^{k-1}$ untuk bilangan prima $p$ dan bilangan bulat
          positif $k$. Contohnya $\varphi(8) = 8 - 4 = 4$, yaitu bilangan $1,3,5,7$.
    \item $\varphi(mn) = \varphi(m)\,\varphi(n)$ bila $\text{PBB}(m,n) = 1$. Karena
          itulah, untuk $n = pq$ hasil kali dua bilangan prima berbeda berlaku
          \[ \varphi(n) = (p-1)(q-1). \]
\end{itemize}
Sifat terakhir adalah inti pembangkitan kunci RSA. Perhatikan bahwa untuk menghitung
$\varphi(n)$ bagi $n = pq$, kita hanya perlu mengetahui $p$ dan $q$ --- bukan
memfaktorkan $n$ dari nol. Sebaliknya, siapa pun yang dapat menghitung $\varphi(n)$
dapat memfaktorkan $n$, sehingga kerahasiaan $\varphi(n)$ setara dengan kerahasiaan
faktor-faktornya.

\subsection{Teorema Euler dan Teorema Fermat Kecil}

Teorema berikut menghubungkan eksponensiasi modular dengan fungsi totient, dan
merupakan mesin di balik kebenaran algoritma RSA.

\begin{theorem}[Teorema Euler]
Jika $a$ dan $n$ relatif prima, yaitu $\text{PBB}(a,n) = 1$, maka
\begin{equation}
    a^{\varphi(n)} \equiv 1 \pmod{n}.
    \label{eq:euler}
\end{equation}
\end{theorem}

Sebagai pemeriksaan, ambil $a = 3$ dan $n = 10$. Karena $\varphi(10) = 4$ dan
$3^4 = 81 = 8 \cdot 10 + 1$, benar bahwa $3^4 \equiv 1 \pmod{10}$.

\begin{theorem}[Teorema Fermat Kecil]
Jika $p$ adalah bilangan prima dan $\text{PBB}(a,p) = 1$, maka
\[ a^{p-1} \equiv 1 \pmod{p}. \]
\end{theorem}

Teorema Fermat Kecil adalah kasus khusus Teorema Euler, diperoleh dengan
mensubstitusi $\varphi(p) = p - 1$. Contohnya, untuk $p = 11$ dan $a = 2$ berlaku
$2^{10} = 1024 = 93 \cdot 11 + 1$, sehingga $2^{10} \equiv 1 \pmod{11}$. Teorema ini
juga menjadi dasar uji keprimaan Fermat dan uji Miller--Rabin yang telah disinggung
sebelumnya.

Perhatikan bahwa Teorema Euler mensyaratkan $\text{PBB}(a,n)=1$. Pada RSA, syarat
ini dipenuhi untuk hampir seluruh pesan; kasus $m$ yang tidak relatif prima terhadap
$n$ hanya terjadi bila $m$ merupakan kelipatan $p$ atau $q$, dan untuk modulus yang
cukup besar peluangnya dapat diabaikan. Penurunan lengkap mengapa
$c^d \equiv m \pmod n$ akan disajikan pada Bab~\ref{chap:rsa}.

\subsection{Akar Primitif dan Logaritma Diskrit}

Misalkan $p$ bilangan prima. Karena setiap $a$ dengan $1 \le a \le p-1$ relatif prima
terhadap $p$, Teorema Fermat Kecil menjamin $a^{p-1} \equiv 1 \pmod p$. Namun
eksponen terkecil yang membuat $a$ kembali ke 1 bisa jauh lebih kecil dari $p-1$.
Jika eksponen terkecil itu tepat $p-1$, maka $a$ disebut \textbf{akar primitif}
(\textit{primitive root}) modulo $p$, dan pangkat-pangkatnya
\[ a^1, a^2, a^3, \dots, a^{p-1} \pmod p \]
menghasilkan seluruh bilangan $1, 2, \dots, p-1$ tepat satu kali masing-masing.

Sebagai contoh, untuk $p = 7$:
\begin{itemize}
    \item $g = 3$ adalah akar primitif, karena pangkat-pangkatnya membentuk siklus
          penuh: $3, 2, 6, 4, 5, 1$ --- keenam anggota $\{1,\dots,6\}$ muncul semua.
    \item $g = 2$ bukan akar primitif, karena $2^1 = 2$, $2^2 = 4$, $2^3 = 8 \equiv 1$,
          sehingga siklusnya hanya berpanjang tiga dan hanya menghasilkan
          $\{1,2,4\}$.
\end{itemize}

Keberadaan akar primitif ini melahirkan persoalan \textbf{logaritma diskrit}: bila
$g$ adalah akar primitif modulo $p$ dan $h$ adalah sembarang anggota $\{1,\dots,p-1\}$,
maka terdapat tepat satu eksponen $x$ dengan
\[ g^x \equiv h \pmod p, \]
dan $x$ disebut logaritma diskrit dari $h$ dengan bilangan pokok $g$. Sebagai contoh,
untuk $p = 41$ dan $g = 7$ diperoleh $7^3 = 343 = 8 \cdot 41 + 15$, sehingga
$7^3 \equiv 15 \pmod{41}$ dan logaritma diskrit dari 15 dengan bilangan pokok 7
adalah $x = 3$.

Menghitung $g^x \bmod p$ dari nilai $x$ yang diketahui mudah dilakukan --- cukup
dengan eksponensiasi modular berulang. Sebaliknya, menemukan $x$ dari $g$ dan $h$
dianggap sulit: tidak ada algoritma klasik yang menyelesaikannya dalam waktu
polinomial terhadap banyaknya digit $p$, dan algoritma terbaik yang dikenal bersifat
subeksponensial. Asimetri inilah yang menjadi dasar protokol Diffie--Hellman pada
Bab~\ref{chap:diffie-hellman} dan algoritma ElGamal pada Bab~\ref{chap:elgamal},
sebagaimana kesulitan pemfaktoran menjadi dasar RSA pada Bab~\ref{chap:rsa}
\citep{menezes1996}.
